\section{Introducción}

Mientras el capítulo de Antecedentes compara logSguarDian contra trabajos previos para identificar los vacíos metodológicos que este proyecto atiende, el presente capítulo desarrolla los conceptos que un lector necesita para entender por qué esos vacíos existen y por qué la arquitectura propuesta los cierra de la forma en que lo hace. Los fundamentos se organizan en el orden en que se articulan dentro del sistema: primero las categorías de ataque que logSguarDian debe reconocer y el entorno de ejecución --Node.js, Express y el modelo de middleware-- en el que opera; después el ecosistema de dependencias de terceros y el contexto local que motivan su necesidad; luego los algoritmos de aprendizaje automático evaluados y seleccionados para el modelo híbrido; y finalmente el proceso de ingeniería de características, las métricas con las que se evalúa el sistema, y las restricciones de serialización y concurrencia bajo las cuales debe operar en producción. Cada eje se cierra relacionando el concepto con una decisión de diseño concreta del proyecto, documentada con mayor detalle técnico en los capítulos de Alcance y Objetivos.

\section{Panorama de amenazas en la capa de aplicación}

El OWASP Top~10 clasifica la inyección como la categoría de riesgo más consistentemente crítica en aplicaciones web desde su primera edición, y la define como cualquier mecanismo por el cual datos no confiables suministrados por un atacante alteran la interpretación de un comando o consulta por parte del intérprete que la recibe \parencite{owasp2021}. Las cuatro categorías cubiertas por logSguarDian son variantes concretas de este mecanismo general, cada una dirigida a un intérprete distinto dentro de la pila de una aplicación Node.js típica.

La inyección SQL (SQLi) ocurre cuando una entrada de usuario se concatena sin sanitizar dentro de una consulta destinada a un motor de bases de datos relacional, permitiendo que el atacante inserte cláusulas SQL adicionales --por ejemplo, un operador \texttt{OR 1=1} o un comentario \texttt{--} que trunca el resto de la consulta original-- y así alterar su lógica de autenticación o extracción de datos. El Cross-Site Scripting (XSS) sigue el mismo principio pero dirigido al intérprete del navegador: un payload que contiene marcado HTML o JavaScript --típicamente una etiqueta \texttt{<script>} o un manejador de eventos como \texttt{onerror}-- se refleja o persiste sin escapar en una respuesta, ejecutándose en el contexto de sesión de otro usuario. El Path Traversal, cuya forma más severa es la Inclusión Local de Archivos (LFI), explota la resolución de rutas del sistema de archivos: una secuencia de ascenso de directorio (\texttt{../}) inyectada en un parámetro que la aplicación concatena directamente a una ruta base permite acceder a archivos fuera del directorio previsto. La inyección de comandos del sistema operativo (CmdI) ocurre cuando una entrada de usuario llega sin sanitizar a una llamada que invoca un shell del sistema operativo, permitiendo el uso de operadores de encadenamiento (\texttt{;}, \texttt{\&\&}, \texttt{|}) o de sustitución de comandos (\texttt{\$(...)}) para ejecutar comandos arbitrarios con los privilegios del proceso de la aplicación.

Las cuatro categorías comparten una propiedad estructural que fundamenta directamente el diseño del pipeline de características de logSguarDian: la evidencia del ataque reside íntegramente dentro de una única petición HTTP --su método, ruta, encabezados, parámetros de consulta y cuerpo-- sin requerir estado acumulado de peticiones anteriores. Esta propiedad es la razón técnica, documentada en el capítulo de Alcance, por la cual el vector de características de logSguarDian opera petición a petición y por la cual ataques de naturaleza inter-request como la denegación de servicio o la fuerza bruta quedan fuera del alcance del proyecto.

\section{Node.js, Express y el modelo event-driven}

Node.js ejecuta el código de la aplicación --incluido el middleware que la protege-- sobre un único hilo de Event Loop, que atiende operaciones de entrada/salida de forma no bloqueante mediante \textit{libuv}, pero que no puede ejecutar dos fragmentos de código JavaScript de forma simultánea dentro de ese hilo \parencite{nodejs2024}. Este modelo, conocido como orientado a eventos (\textit{event-driven}), sustituye el patrón clásico de un hilo del sistema operativo por conexión con un bucle único que procesa callbacks a medida que sus operaciones de E/S asociadas se resuelven, lo cual permite atender miles de conexiones concurrentes con un consumo de memoria muy inferior al de un servidor multihilo tradicional, a costa de que cualquier código síncrono de larga duración bloquee la atención de todas las demás conexiones mientras se ejecuta.

Express se construye sobre este modelo como un framework de enrutamiento que organiza el procesamiento de una petición HTTP en una cadena de funciones \textit{middleware}, cada una con acceso al objeto de petición (\texttt{req}) y de respuesta (\texttt{res}), y con la capacidad de delegar el control a la siguiente función de la cadena mediante una llamada explícita (\texttt{next()}) o de interrumpirla enviando una respuesta. Esta arquitectura de cadena de middleware es la que permite insertar un componente de seguridad como logSguarDian delante de la lógica de negocio de la aplicación sin modificar su código: basta con registrarlo como una función más en la cadena, ejecutada antes que las rutas de la aplicación anfitriona. La consecuencia directa del modelo de un solo hilo para el diseño de logSguarDian es que cualquier operación de cómputo intensivo ejecutada de forma síncrona dentro de ese middleware --como la inferencia de un modelo de aprendizaje automático-- bloquearía el Event Loop y, con él, todas las peticiones concurrentes de la aplicación anfitriona, sin importar cuán rápida sea esa operación en términos absolutos; esta restricción se retoma en la sección de Edge AI de este capítulo, donde se describe el mecanismo con el que logSguarDian la resuelve.

\section{Middleware y Runtime Application Self-Protection (RASP)}

La posición de logSguarDian dentro de la cadena de middleware de Express no es solo un detalle de integración: define su categoría arquitectónica frente a otros mecanismos de defensa de la capa de aplicación. \textcite{yuan2012} formalizan, dentro de una taxonomía más amplia de sistemas de software autoprotector, la categoría de autoprotección de aplicaciones en tiempo de ejecución (RASP, por sus siglas en inglés) como aquella que instrumenta directamente el proceso de la aplicación para detectar y mitigar amenazas con visibilidad de su estado de ejecución interno, en contraposición a los mecanismos perimetrales que observan únicamente el tráfico de red que entra y sale del sistema. logSguarDian se posiciona formalmente dentro de esta categoría: al ejecutarse como middleware dentro del mismo proceso Node.js que la aplicación protegida, tiene acceso directo al objeto \texttt{req} de Express ya parseado por el framework, sin necesidad de reimplementar la normalización de protocolo HTTP que un mecanismo externo debe resolver de forma genérica y, por definición, con menor información de contexto que la propia aplicación.

La arquitectura con la que históricamente se ha contrastado un RASP es el cortafuegos de aplicación web (WAF) de arquitectura perimetral, ejemplificado en este proyecto por ModSecurity operando con el OWASP Core Rule Set (CRS). Un WAF de este tipo se despliega como proxy inverso frente al servidor de aplicaciones e inspecciona cada petición mediante un conjunto de reglas de coincidencia de patrones y funciones de normalización, entre ellas la decodificación de URL, la remoción de espacios redundantes y la detección de sintaxis SQL mediante análisis léxico (\textit{libinjection}), antes de reenviarla o bloquearla \parencite{owaspcrs2024}. Su fortaleza es la cobertura inmediata de firmas de ataque conocidas sin requerir instrumentación de la aplicación protegida; su limitación estructural es que opera sin visibilidad del contexto interno de la aplicación --no conoce el modelo de datos, la lógica de negocio ni el resultado real de procesar la petición-- y que su eficacia depende de que el atacante no logre evadir las reglas de normalización y coincidencia mediante técnicas de ofuscación del payload.

La Tabla \ref{cuadro:waf_vs_rasp} contrasta ambas arquitecturas en las dimensiones relevantes para la decisión de diseño de este proyecto y para la interpretación de la Configuración~3 de la evaluación, en la cual ambas arquitecturas operan en capas sobre el mismo tráfico.

\begin{table}[h]
\small
\begin{tabular}{|p{2.6cm}|p{5.3cm}|p{5.3cm}|}
\hline
\textbf{Dimensión} & \textbf{WAF perimetral (ModSecurity/CRS)} & \textbf{RASP embebido (logSguarDian)} \\ \hline
Ubicación & Proxy inverso, proceso independiente & Middleware dentro del proceso de la aplicación \\ \hline
Visibilidad de contexto & Solo tráfico de red entrante/saliente & Objeto de petición ya parseado por el framework \\ \hline
Mecanismo de detección & Coincidencia de patrones y firmas sobre payload normalizado & Clasificación estadística sobre un vector de características extraído de la petición \\ \hline
Cobertura frente a variantes no vistas & Limitada a lo cubierto por las reglas de firma & Provista por el componente no supervisado del modelo híbrido \\ \hline
Dependencia de infraestructura externa & Requiere proxy y mantenimiento de reglas & Ninguna; opera como dependencia de la aplicación \\ \hline
\end{tabular}
\caption[Comparación entre arquitectura WAF perimetral y RASP embebida]{Comparación entre arquitectura WAF perimetral y RASP embebida en las dimensiones relevantes para este proyecto. Elaboración propia a partir de \textcite{yuan2012} y \textcite{owaspcrs2024}.}
\label{cuadro:waf_vs_rasp}
\end{table}

\section{Ecosistema npm y seguridad en la cadena de suministro}

logSguarDian se distribuye como un paquete del registro npm, la misma vía por la que la aplicación anfitriona obtiene la inmensa mayoría de sus dependencias. \textcite{zimmermann2019} documentan, mediante el análisis sistemático del grafo de dependencias de npm, que un paquete promedio confía implícitamente en 79 paquetes de terceros y 39 mantenedores distintos de forma transitiva, y que la ausencia de un proceso de revisión centralizado antes de la publicación permite que un único paquete comprometido --ya sea por una cuenta de mantenedor secuestrada o por código malicioso introducido deliberadamente-- se propague a decenas de miles de proyectos que dependen de él directa o indirectamente. Este hallazgo tiene dos implicaciones directas para el diseño de logSguarDian.

La primera es una restricción de diseño: una librería cuyo propósito es proteger una aplicación no puede, sin contradecir ese propósito, introducir ella misma una superficie de ataque de cadena de suministro significativa. Esto motiva mantener el número de dependencias de tiempo de ejecución de logSguarDian deliberadamente acotado --\texttt{onnxruntime-node} para la inferencia, \texttt{better-sqlite3} para el almacenamiento local de eventos, y las dependencias del propio framework Express-- en lugar de adoptar un ecosistema extenso de utilidades de terceros para funciones que pueden resolverse con la biblioteca estándar de Node.js. La segunda implicación es de alcance: la superficie de riesgo documentada por \textcite{zimmermann2019} corresponde a paquetes que ejecutan código arbitrario durante su instalación o su uso, un vector estructuralmente distinto de las cuatro categorías de inyección en tráfico HTTP cubiertas por este proyecto, razón por la cual la seguridad de la cadena de suministro de npm se documenta aquí como parte del contexto que motiva un diseño de dependencias mínimas, pero no como una categoría de ataque que el modelo de clasificación de logSguarDian esté diseñado para detectar.

\section{Panorama de ciberseguridad en Guatemala y el ecosistema tech local}

La motivación de este proyecto, desarrollada con mayor detalle en el capítulo de Justificación, se apoya en una condición estructural del contexto en el que se ubica: un ecosistema de startups tecnológicas que opera predominantemente sin equipo de seguridad dedicado, en una región donde la exposición a ciberataques ha crecido de forma pronunciada en los años recientes. El informe de \textcite{mastercard2024} documenta un incremento del 200\% en ciberataques registrados en Guatemala entre 2023 y 2024, y ubica al país entre los diez de mayor exposición cibernética de Latinoamérica. Esta cifra no es, en sí misma, un antecedente técnico del diseño de logSguarDian --no informa decisiones de arquitectura, algoritmo o umbral-- pero sí fundamenta la premisa sobre la cual se justifica la existencia del proyecto: que la brecha entre la superficie de ataque documentada por estándares como OWASP Top~10 \parencite{owasp2021} y la capacidad real de organizaciones pequeñas para mitigarla no es un escenario hipotético, sino una condición medible del entorno en el que logSguarDian está pensado para operar. Este panorama se retoma en el capítulo de Justificación para argumentar la relevancia social y económica del proyecto, y no se desarrolla aquí más allá de establecer el marco contextual en el que se inserta la motivación técnica de las secciones siguientes.

\section{Aprendizaje automático supervisado y no supervisado en ciberseguridad}

El modelo híbrido de logSguarDian combina un clasificador supervisado y un detector de anomalías no supervisado, decisión que se fundamenta en la distinción formal establecida por \textcite{chandola2009}: la detección supervisada requiere ejemplos etiquetados de cada clase de anomalía que se desea reconocer y ofrece alta precisión sobre esas clases específicas, mientras que la detección no supervisada de una sola clase modela únicamente la distribución de lo normal y señala como anómalo cualquier punto que se aleje de ella, sin necesidad de haber observado ejemplos de la anomalía durante el entrenamiento. Ambos paradigmas son casos particulares de los métodos de ensamble, que \textcite{dietterich2000} caracteriza como algoritmos que combinan las predicciones de múltiples modelos base para obtener un clasificador agregado cuyo error es, bajo condiciones razonables de diversidad entre los modelos base, menor que el de cualquiera de ellos por separado. Esta distinción formal es la que permite justificar, más adelante en este capítulo, por qué ninguno de los dos paradigmas por sí solo --supervisado o no supervisado-- es suficiente como criterio único de bloqueo, y por qué logSguarDian los combina asignándoles roles distintos: el componente supervisado como única autoridad de bloqueo sobre las categorías de ataque conocidas, y el componente no supervisado como una capa complementaria de enriquecimiento de registro frente a variantes no vistas durante el entrenamiento.

\subsection{Random Forest}

El componente supervisado es un Random Forest, introducido por \textcite{breiman2001} como un método que construye un ensamble de árboles de decisión mediante dos fuentes independientes de aleatoriedad: cada árbol se entrena sobre una muestra bootstrap distinta del conjunto de entrenamiento, y en cada nodo de división se evalúa únicamente un subconjunto aleatorio de las características disponibles en lugar de todas ellas. La predicción final se obtiene por voto mayoritario entre los árboles. Esta doble aleatorización reduce la correlación entre los árboles individuales y, con ella, la varianza del ensamble frente a un único árbol de decisión, sin requerir normalización previa de las características ni supuestos sobre su distribución. Estas dos propiedades --tolerancia a características de escalas heterogéneas y complejidad de inferencia lineal en el número de árboles-- son la justificación técnica directa, documentada también en el capítulo de Antecedentes, de por qué Random Forest fue seleccionado frente a alternativas de mayor capacidad como las redes neuronales profundas para un componente que debe ejecutarse dentro del presupuesto de latencia de un middleware.

\subsection{Isolation Forest}

El componente no supervisado es un Isolation Forest, propuesto por \textcite{liu2008} bajo un principio distinto al de la mayoría de los detectores de anomalías de la época, que suelen modelar la densidad o la distancia entre puntos: en lugar de perfilar qué aspecto tiene un punto normal, Isolation Forest mide qué tan fácil es aislar un punto del resto mediante particiones aleatorias sucesivas del espacio de características. Una observación anómala, por estar más alejada de la masa de datos normales, requiere en promedio un número menor de particiones para quedar aislada en su propia hoja que una observación típica; esa longitud de camino promedio, normalizada, constituye el puntaje de anomalía. Esta propiedad hace que el algoritmo tenga la misma complejidad computacional lineal que Random Forest sin requerir el cálculo de distancias par a par que sí exigen otros métodos de detección de anomalías, lo cual lo hace compatible con las mismas restricciones de inferencia embebida que motivaron la selección del componente supervisado.

\section{Redes neuronales ligeras (MLP) y autoencoders}

Antes de fijar Random Forest e Isolation Forest como los algoritmos del modelo híbrido, el diseño de logSguarDian consideró dos alternativas basadas en redes neuronales, ambas descritas formalmente por \textcite{goodfellow2016}. Un perceptrón multicapa (MLP) es una red neuronal compuesta por capas densas sucesivas de neuronas con una función de activación no lineal, entrenada mediante descenso de gradiente y retropropagación del error para aprender fronteras de decisión que un modelo lineal no podría representar; aplicado como componente supervisado, ofrecería una capacidad de representación mayor que un ensamble de árboles a costa de un proceso de entrenamiento más sensible a la escala de las características y de una arquitectura de inferencia basada en multiplicaciones de matrices en cada capa, cuyo costo computacional crece con la profundidad y el ancho de la red. Un autoencoder es una arquitectura no supervisada compuesta por un codificador que comprime la entrada a una representación de menor dimensión y un decodificador que intenta reconstruirla; entrenado únicamente sobre tráfico benigno, el error de reconstrucción sobre una entrada nueva sirve como puntaje de anomalía, dado que un error alto indica una entrada que la red no aprendió a representar bien porque difiere de la distribución de entrenamiento, de forma conceptualmente análoga al puntaje de aislamiento de Isolation Forest.

Ambas alternativas fueron consideradas y descartadas a partir de la lógica documentada en el capítulo de Antecedentes a partir de la revisión de \textcite{berman2019}, y la medición empírica de este proyecto sobre el mismo corpus y contrato de características usados en producción confirma esa lógica de forma desigual entre los dos componentes. Para el componente supervisado, Random Forest (macro F1=0.9776) iguala en la práctica a un MLP equivalente (macro F1=0.9719) con una latencia de inferencia menor (0.0062~ms frente a 0.0074~ms), de modo que la ganancia de precisión que ofrecería la red neuronal es marginal frente al costo adicional, mínimo pero medible, que introduce.

Para el componente no supervisado, la misma medición matiza la justificación original. Un autoencoder entrenado sobre el contrato de características de Isolation Forest iguala su recall y su tasa de falsos positivos (0.9158/0.0517 frente a 0.9125/0.0546) con una latencia de inferencia y un tamaño de modelo serializado órdenes de magnitud menores (0.0065~ms frente a 0.7441~ms; 19.3~KB frente a 2.06~MB). Este costo diferencial no proviene de que Isolation Forest sea, como familia de algoritmo, más costosa que una red neuronal equivalente, sino de la configuración adoptada en producción --200 árboles con \texttt{max\_samples=4096}, fijados por razones de estabilidad de calibración-- y de que su exportación a ONNX resulta, en la práctica, un subgrafo menos eficiente que el de Random Forest. La decisión de mantener Isolation Forest se sostiene, entonces, por una razón arquitectónica distinta a la de costo bruto de inferencia: al no tener autoridad de bloqueo sobre la respuesta al cliente y ejecutarse de forma asíncrona en el \textit{worker pool} dedicado documentado en el capítulo de Alcance, su latencia queda absorbida fuera del camino crítico de la petición, lo que hace que la comparación de latencia bruta sea menos determinante para este componente que para el supervisado. La consideración explícita de ambas alternativas, y no solo la ausencia de su implementación, es lo que permite documentar la selección de Random Forest e Isolation Forest como una decisión de diseño evaluada --y en el caso de Isolation Forest, sostenida por arquitectura y no por costo bruto-- y no como la única opción contemplada.

\section{Ingeniería de características para datos de seguridad HTTP}

Ningún algoritmo de aprendizaje automático opera directamente sobre un objeto de petición HTTP: requiere que ese objeto se transforme primero en un vector numérico de dimensión fija, proceso que \textcite{zheng2018} denominan ingeniería de características y describen como la etapa que determina en la práctica gran parte del desempeño final de un modelo, con independencia del algoritmo de clasificación elegido sobre esas características. Entre las técnicas que documentan y que logSguarDian adopta de forma directa se encuentran el conteo de ocurrencias de patrones específicos dentro de una cadena de texto (por ejemplo, el número de operadores lógicos SQL o de secuencias de ascenso de directorio dentro de un payload), la codificación de atributos categóricos --como el método HTTP o la presencia de un encabezado particular-- mediante variables indicadoras, y la derivación de estadísticos simples sobre texto libre, como la longitud de una cadena o la proporción de caracteres no alfanuméricos que contiene.

El vector de 75 características de logSguarDian, descrito en detalle en el capítulo de Alcance y en la especificación de características del proyecto, aplica estos principios generales a seis grupos funcionales orientados a capturar los marcadores semánticos de cada una de las cuatro categorías de ataque cubiertas, además de características estadísticas transversales y metadatos de la petición. La adaptación específica de este proyecto frente al cuerpo general de técnicas documentado por \textcite{zheng2018} es la selección del objeto \texttt{req} de Express --y no un flujo de red o un documento de texto genérico-- como fuente de esas características, decisión motivada por el vacío metodológico identificado en el capítulo de Antecedentes respecto a la ausencia de literatura de referencia sobre ingeniería de características a nivel de petición HTTP individual.

\section{Detección de anomalías y ataques de día cero}

Un ataque de día cero, en el sentido más general en que el término se emplea en este trabajo --y no en su acepción más estricta de vulnerabilidad de software desconocida por el fabricante--, es una variante de payload que no se parece a ningún ejemplo visto durante el entrenamiento del componente supervisado: una técnica de ofuscación nueva, una combinación de operadores no representada en el corpus, o una categoría de ataque que ni siquiera fue anticipada como clase durante el diseño del modelo. Un clasificador puramente supervisado como Random Forest no tiene, por construcción, forma de reconocer una entrada que no se parezca a ninguna de las clases sobre las que fue entrenado: la asignará a la clase más cercana en el espacio de características, con una confianza que no refleja necesariamente qué tan atípica es esa entrada frente a la distribución de entrenamiento.

Esta limitación es la que motiva, siguiendo la distinción de \textcite{chandola2009} desarrollada en la sección anterior, la inclusión de Isolation Forest como una capa de detección complementaria que no depende de haber visto ejemplos previos de una técnica de ataque para señalarla como estadísticamente atípica. En logSguarDian, esta capacidad se documenta operando en dos escenarios distintos: como diagnóstico en tiempo real dentro del middleware --donde el puntaje de anomalía enriquece el registro de una petición sin constituir por sí solo un criterio de bloqueo-- y, de forma más directa, dentro del pipeline de reentrenamiento continuo del proyecto, donde la combinación de Isolation Forest con un algoritmo de agrupamiento por densidad permitió aislar, sin etiqueta previa, un payload poliglota que combinaba simultáneamente patrones de las cuatro categorías de ataque cubiertas en una única petición, algo que el componente supervisado por sí solo clasificó erróneamente como una sola categoría. Este resultado ilustra tanto el valor como el límite de la detección no supervisada de día cero: identifica como atípico lo que se desvía lo suficiente de la distribución de entrenamiento en varios ejes a la vez, pero no ofrece garantía de que cualquier técnica nueva individual, por sí sola, cruce ese umbral de atipicidad.

\section{Métricas de evaluación: F1-score, AUC-ROC y matriz de confusión}

La evaluación de un clasificador de amenazas no puede apoyarse únicamente en la exactitud (\textit{accuracy}), la proporción de predicciones correctas sobre el total, porque un corpus con clases fuertemente desbalanceadas --como es el caso de tráfico benigno frente a tráfico malicioso en un entorno realista-- permite que un clasificador que etiquete todo como benigno obtenga una exactitud alta sin utilidad práctica alguna. Por esta razón, los Objetivos específicos de este proyecto se formulan en términos de precisión, recuperación (\textit{recall}) y su combinación en el F1-score, calculadas a partir de la matriz de confusión de verdaderos positivos ($VP$), falsos positivos ($FP$), verdaderos negativos ($VN$) y falsos negativos ($FN$):

\begin{equation}
\text{Precisión} = \frac{VP}{VP + FP} \qquad
\text{Recall} = \frac{VP}{VP + FN} \qquad
\text{F1} = 2 \cdot \frac{\text{Precisión} \cdot \text{Recall}}{\text{Precisión} + \text{Recall}}
\end{equation}

La precisión mide qué proporción de las peticiones bloqueadas por el modelo corresponden efectivamente a un ataque, y su complemento (uno menos la precisión, ponderado por la tasa de falsos positivos) es la métrica directamente relevante para el criterio de aceptación del componente no supervisado, dado que un Isolation Forest con una tasa de falsos positivos alta bloquearía tráfico legítimo con la misma frecuencia con la que detecta anomalías reales. El recall mide qué proporción de los ataques reales el modelo efectivamente detecta, y es la métrica directamente relevante para el criterio de cobertura del 80\% por categoría de ataque establecido en el tercer objetivo específico. El F1-score, al ser la media armónica de ambas, penaliza a un clasificador que sacrifique una métrica para maximizar la otra --por ejemplo, bloqueando indiscriminadamente para maximizar recall a costa de la precisión-- y es, por esta razón, la métrica elegida como criterio de aceptación del componente supervisado en el segundo objetivo específico. \textcite{fawcett2006} formaliza además el análisis de curvas ROC (\textit{Receiver Operating Characteristic}) y de su área bajo la curva (AUC) como herramienta complementaria para evaluar el compromiso entre tasa de verdaderos positivos y tasa de falsos positivos a través de distintos umbrales de decisión, principio aplicado en este proyecto durante la calibración del umbral de bloqueo del componente supervisado y del umbral de anomalía del componente no supervisado.

\section{Edge AI e inferencia en tiempo de ejecución con ONNX}

La ejecución de modelos de aprendizaje automático dentro de un proceso con restricciones estrictas de latencia y memoria, sin comunicación con servicios externos, está documentada de forma consolidada en la literatura de Edge AI, cuyo texto de referencia es el de \textcite{warden2019}. Los autores establecen los principios de diseño para desplegar modelos ligeros bajo presupuestos de memoria y tiempo de inferencia del orden de milisegundos, incluyendo técnicas de cuantización y formatos de serialización optimizados; esta disciplina, desarrollada casi en su totalidad sobre microcontroladores y dispositivos embebidos de hardware dedicado, provee el principio general que logSguarDian traslada a un dominio de ejecución distinto: el de un único proceso Node.js de un servidor de aplicaciones.

El entrenamiento de los modelos de logSguarDian se realiza en Python sobre \textit{scikit-learn}, mientras que su inferencia en producción ocurre dentro de un proceso Node.js, dos entornos de ejecución sin compatibilidad binaria directa entre sí. Open Neural Network Exchange (ONNX) resuelve esta discontinuidad como un formato abierto de representación de modelos de aprendizaje automático, que define un grafo de cómputo y un conjunto de operadores estándar independientes del framework en el que el modelo fue entrenado \parencite{bai2019}, permitiendo que un modelo serializado en ONNX se ejecute mediante un motor de inferencia distinto (en este caso, \textit{onnxruntime-node}) al que se usó para entrenarlo. Esta propiedad es la que hace viable, en primer lugar, entrenar los modelos híbridos de logSguarDian con el ecosistema maduro de \textit{scikit-learn} en Python y, en segundo lugar, distribuirlos embebidos dentro de un paquete npm sin requerir que el proceso de inferencia dependa de un intérprete de Python en producción. La contrapartida de esta traducción entre representaciones es la necesidad de verificar que ambas produzcan predicciones equivalentes sobre las mismas entradas --verificación de paridad que este proyecto ejecuta de forma sistemática tras cada cambio de modelo, dado que una divergencia entre el modelo evaluado en Python y el modelo efectivamente servido en Node.js invalidaría las métricas reportadas en este trabajo.

Resta resolver, sin embargo, la restricción de concurrencia descrita en la sección de Node.js y Express de este capítulo: un motor de inferencia ONNX ejecutado de forma síncrona en el hilo principal bloquearía el Event Loop durante cada predicción. La documentación oficial de Node.js introduce \texttt{worker\_threads} precisamente como el mecanismo para delegar trabajo de cómputo intensivo a hilos del sistema operativo independientes del Event Loop, comunicándose con el hilo principal mediante paso de mensajes en lugar de memoria compartida \parencite{nodejs2024}. Este mecanismo es el fundamento técnico directo del primer objetivo específico del proyecto, que exige que el pipeline de ingeniería de características y la inferencia del modelo híbrido se ejecuten de forma asíncrona mediante \texttt{worker\_threads} sin introducir una latencia adicional superior a 5~ms por petición en el servidor anfitrión, y de la arquitectura de \textit{worker pool} dedicado documentada en el capítulo de Alcance. Es también la razón por la cual la librería delimita su compatibilidad a partir de la versión 12 de Node.js, la primera en la que \texttt{worker\_threads} está disponible de forma estable sin necesidad de una bandera experimental.

\section{Construcción y etiquetado de datasets de tráfico HTTP}

Entrenar un clasificador supervisado como Random Forest requiere un corpus en el que cada petición esté etiquetada con su categoría real, condición que logSguarDian resuelve mayoritariamente aprovechando la etiqueta ya asignada por la publicación original de cada dataset público que incorpora --SR-BH~2020/CAPEC, modsec-learn, registros de honeypot de ModSecurity, SecLists, PayloadsAllTheThings y varios conjuntos de Kaggle, entre otros, con licencia y procedencia verificadas-- sin necesidad de anotación manual adicional. Esta estrategia deja, sin embargo, un problema distinto sin resolver: ningún dataset público disponible cubre de antemano los vacíos de cobertura específicos que solo se hacen evidentes al validar el modelo ya entrenado --por ejemplo, un patrón de navegación benigna subrepresentado que un análisis de exclusión por fuente (LOSO) revela como punto débil, o una técnica de inyección de comandos cuya confianza de detección una prueba de ablación muestra insuficiente--. Para ese problema, distinto al de la etiquetación del corpus público, logSguarDian recurre a un suplemento acotado de tráfico sintético generado deliberadamente para cubrir cada vacío identificado.

\textcite{shiravi2012} proponen el enfoque adoptado como base metodológica para ese suplemento: en lugar de etiquetar tráfico observado después de generarlo, se despliega un entorno de prueba controlado, se simula tráfico legítimo representativo y se ejecutan herramientas de ataque automatizadas de forma deliberada, de modo que la etiqueta de cada muestra se conoce de antemano por construcción --\textit{ground truth by design}-- y no por inferencia posterior. Este principio evita la necesidad de anotación manual y garantiza una correspondencia exacta entre la muestra y su etiqueta, a costa de que la representatividad de esa porción del corpus dependa enteramente de qué tan bien el entorno de prueba y las herramientas de generación aproximan el tráfico real que el sistema enfrentará en producción, supuesto que el capítulo de Alcance documenta explícitamente como una condición verificada mediante análisis exploratorio del corpus resultante.

logSguarDian aplica este mismo principio metodológico --despliegue de un entorno de prueba controlado mediante Docker Compose, simulación de tráfico benigno y ejecución deliberada de payloads de ataque-- únicamente al suplemento sintético, dirigiéndolo además a un objeto distinto del que documenta la literatura de referencia sobre generación de datasets de intrusión: en lugar de capturar metadatos de flujo de red (duración de conexión, bytes transferidos, número de paquetes), como es el caso de los datasets de la familia CICIDS derivados de \textcite{shiravi2012}, cada muestra sintética del corpus de logSguarDian es un vector de características extraído del objeto \texttt{req} de Express mediante el mismo pipeline de ingeniería de características que se ejecuta en producción, descrito en la sección correspondiente de este capítulo, y ese mismo pipeline se aplica después sobre las muestras de origen público para producir un vector de entrada homogéneo con independencia de la procedencia de cada fila. Esta adaptación es la que permite que el corpus de entrenamiento sea representativo de la semántica de una petición HTTP individual y no únicamente de patrones de tráfico a nivel de conexión.

\section{Síntesis: articulación de los fundamentos en la arquitectura de logSguarDian}

Los fundamentos desarrollados en este capítulo no son piezas de conocimiento independientes: se articulan directamente en las decisiones de diseño de logSguarDian. La caracterización de las cuatro categorías de ataque como fenómenos intra-request delimita la unidad de análisis del sistema a una petición HTTP individual, y el modelo de middleware de Express sobre el Event Loop de un solo hilo de Node.js define tanto el punto de integración de logSguarDian como su restricción de latencia más estricta. La distinción entre arquitectura perimetral y arquitectura RASP embebida fundamenta la decisión de operar dentro del proceso de la aplicación en lugar de como proxy independiente, mientras que la seguridad de la cadena de suministro de npm y el panorama de ciberseguridad del contexto local delimitan, respectivamente, una restricción de diseño sobre las dependencias del proyecto y la motivación de su existencia. La complementariedad entre detección supervisada y no supervisada fundamenta la combinación de Random Forest e Isolation Forest en un modelo híbrido --evaluada frente a las alternativas de redes neuronales descritas en este capítulo, y sostenida en un caso por su costo de inferencia y en el otro por su arquitectura de ejecución asíncrona-- y provee el mecanismo conceptual de detección de variantes de día cero. Los principios generales de ingeniería de características se adaptan al objeto de petición de Express para producir el vector de entrada de ambos modelos, entrenados sobre un corpus compuesto predominantemente por datasets públicos con etiqueta propia y complementado con un suplemento sintético construido bajo el principio de etiquetado determinista documentado en la literatura de generación de datasets de intrusión, y evaluados con las métricas de precisión, recall y F1-score que expresan formalmente los criterios de aceptación de los objetivos específicos de este trabajo. Finalmente, ONNX y \texttt{worker\_threads} resuelven, respectivamente, la discontinuidad entre el entorno de entrenamiento en Python y el entorno de inferencia en Node.js, y la restricción de concurrencia de un solo hilo bajo la cual ese entorno de inferencia debe operar sin degradar el servicio que protege.
